\section{Scaling and coupling estimates}
\label{subsec:rgnl-scaled-correction}

Let $S_L$ be the field scaling of
\exteq{30}{NormalizedContraction}. The nonlinear correction is
\begin{equation}\label{eq:rgnl-scaled-correction}
 E_k^*(m,\mathscr E)(Y;\mathscr A,\mathscr C)
  =\sum_{\bar X^L=LY}S_L
        E_k^\#(m,\mathscr E)(X;\mathscr A_L,\mathscr C_L).
\end{equation}
This formula uses field scaling and summation over fine
polymers. Quadratic adjustment enters the incoming
activity term in the full step below.
The domain inclusions
\exteq{29}{NormalizedContraction} ensure
$(\mathscr A_L,\mathscr C_L)
\in\widetilde R_k^-\times\mathfrak X_k$.

\begin{proposition}\label{prop:rgnl-scaled-correction}
For $(m,\mathscr E)\in\operatorname{Ball}_k^x(4)$,
\[
 E_k^*(m,\mathscr E)\in K_{k+1}^{x,\C},\qquad
                  \|E_k^*(m,\mathscr E)\|_{k+1}^x
                                      \le e^{1/4-6\varepsilon}.
\]
For real data its value belongs to $K_{k+1}^x$.
On any affine complex line in the data space its value
is norm continuous on the part in the closed ball and
norm analytic in its interior.
For two data in $\operatorname{Ball}_k^x(1)$,
\begin{equation}\label{eq:rgnl-coupling-Lipschitz}
 \|E_k^*(m_1,\mathscr E_1)-E_k^*(m_2,\mathscr E_2)\|_{k+1}^x
 \le e^{1/4-6\varepsilon}
       \left(\frac{|m_1-m_2|}{w_k^{\mathrm m}}
              +\frac{\|\mathscr E_1-\mathscr E_2\|_k^x}
                                                    {w_k^{\mathrm E}}\right).
\end{equation}
The second quotient is absent at $k=0$.
Restriction to $\mathscr C=0$ gives the correction formed
without an external slot from $(m,\mathscr E|_0)$.
\end{proposition}

\begin{proof}
The scaling identity gives
\[
 \|S_LQ\|_{h_{k+1}}
   =\|Q\|_{(L^{-1}h_{k+1,1},L^{-1-\alpha}h_{k+1,2})}
   \le\|Q\|_{h_k/4},
\]
since $h_{k+1,1}=L^{-1/8}h_{k,1}$,
$h_{k+1,2}=L^{-1/8+\varepsilon/2}h_{k,2}$ and $L\ge4$.
The polymer sum in \eqref{eq:rgnl-scaled-correction}
uses $Ld_M(Y)\le d_M(X)$ and $X\subseteq LY$.
Splitting $a_3=(a_3-\kappa_0)+\kappa_0$ therefore yields
\begin{align*}
 \sum_{\bar X^L=LY}\exp(-a_3d_M(X))
  &\le \exp(-L(a_3-\kappa_0)d_M(Y))
                         \sum_{X\subseteq LY}\exp(-\kappa_0d_M(X))\\
  &\le c_XK_0L^3
                  \exp(-[L(\kappa-6\kappa_0-4)-1]d_M(Y))
   \le c_XK_0L^3\exp(-\kappa d_M(Y)).
\end{align*}
Together with \eqref{eq:rgnl-log-bound} this gives
\[
 \|E_k^*\|_{k+1}^x
  \le c_XK_0L^3c_S'H_{\mathrm{cl}}
  =504c_X^5K_0^5\exp(3\kappa)L^3e^{1/4-5\varepsilon}
  \le e^{1/4-6\varepsilon},
\]
using $c_SL^3e^\varepsilon\le1$.

All generators of a summand lie over $X$ or at a pair
admissible for $X$. Scaling maps the ordinary sites
into $Y$; it maps each of the six paths into the
corresponding scaled path and maps
$(x,x^{(3)})$ to $(x/L,(x/L)^{(3)})$.
These pairs are admissible for $Y$, including
intracluster pairs. A bond meeting $X$ scales to
one meeting $Y$, giving local connection dependence.
The scaled-domain inclusion and coefficient analyticity
give analytic kernels.
The map $(S,a)\mapsto(S,a/L)$ identifies
$\mathbb G_k^L$ with $\mathbb G_{k+1}$ and intertwines
its generator actions with $S_L$. For gauge transformations
use $\Lambda_L(c)=L^{1/2}\Lambda(c/L)$ and
$(d\Lambda)_L=d(\Lambda_L)$.
A gauge relation on the unit neighborhood of $Y$
then holds on the unit neighborhood of every
$X\subseteq LY$. Charge conjugation and reality also
commute with the real scaling factors.
\Cref{prop:rgnl-symmetries} thus gives all the
defining symmetries of $K_{k+1}^{x,\C}$, and the real
space assertion follows from its reality clause.

For norm analyticity, take data
$(m_a+t m_b,\mathscr E_a+t\mathscr E_b)$ in the closed
ball and denote their parameter set by $T$.
It is closed and convex. The finite activity $F_t(X)$
is affine in $t$ by \cref{prop:rgnl-finite-bounds}.
The polynomial approximants to $E_k^\#$ are therefore
polynomials in $t$. Reblock them by
\eqref{eq:rgnl-scaled-correction}; the resulting polynomials
$f_N(t)$ converge uniformly on $T$ in the full activity
norm, since both finite polymer sets are fixed and the
error bound in \cref{subsec:rgnl-polynomial-approximation} is uniform.
Their values have bounded full norm.
If $T$ has more than one point, it contains infinitely
many points. The coefficients of each polynomial are
finite numerical linear combinations of its values at
distinct points of $T$, by polynomial interpolation.
Hence these coefficients have finite full norm, and the
uniform limit is norm continuous on $T$.
If $T$ is empty or a singleton there is no continuity
or interior analyticity assertion to check.

For a closed disk $|t-t_0|\le r$ in $T$, take its
coefficientwise Cauchy integrals
\[
 f_j=\frac1{2\pi i}\int_{|t-t_0|=r}
                             f(t)(t-t_0)^{-j-1}\,\dd t.
\]
They have analytic connection kernels, and the support,
symmetry and balance conditions pass through this
complex linear integral. Moreover
$\|f_j\|_{k+1}^x\le e^{1/4-6\varepsilon}r^{-j}$.
Coefficientwise Cauchy's formula, with this summable
norm bound, gives
$f(t)=\sum_{j\ge0}f_j(t-t_0)^j$ in activity norm for
$|t-t_0|<r$. This proves the stated analyticity.

To obtain \eqref{eq:rgnl-coupling-Lipschitz}, vary first
the mass, keeping $\mathscr E_1$ fixed. With
$\Delta m=m_1-m_2\ne0$, the line
$(m_2+t\Delta m,\mathscr E_1)$ remains in the large
ball for $|t|\le r=3w_k^{\mathrm m}/|\Delta m|$.
Here $r\ge3/2$. The contour formula on any circle
$1<r'<r$ gives
\[
 \|f(1)-f(0)\|_{k+1}^x
        \le\frac{e^{1/4-6\varepsilon}}{r'-1}.
\]
Let $r'\uparrow r$ and use $r-1\ge r/3$.
The bound is the first quotient on the right of
\eqref{eq:rgnl-coupling-Lipschitz}.
Vary the activity with mass $m_2$ fixed and radius
$3w_k^{\mathrm E}/\|\mathscr E_1-\mathscr E_2\|_k^x$
to obtain the second quotient.
The zero differences require no contour.

Finally restriction commutes with extraction:
$m(\mathscr E)=m(\mathscr E|_0)$ and
$(R\mathscr E)|_0=R(\mathscr E|_0)$.
Every fermion operator uses the sum of its first
connection and $\mathscr C$; photon operations and
measures do not use $\mathscr C$.
At $\mathscr C=0$ the two constructions therefore have
identical finite expressions, connected series and
scaled corrections. Analyticity without a slot is
understood in the first connection alone.
\end{proof}

The possible H\"older legs can also be read directly.
In $G_2,G_4$ an existing pair generator is changed by
addition of a fluctuation field; this creates no new
pair index. The mass contours contain no pair
generators. The reconstruction substitution and the
first-step vertex introduce only pairs $(x,x^{(3)})$.
Every later sum, product and integral preserves this
set of available indices. Thus the correction has
only scaled incoming pair indices and these scaled
type-$3$ indices. At $k=0$ it has only the latter,
of length $\eta_1$, since the incoming activity is zero.

